\documentclass[11pt, a4paper]{article}
\usepackage[a4paper, top=2.5cm, bottom=2.5cm, left=2cm, right=2cm]{geometry}
\usepackage{amsmath,amssymb,amsthm,microtype}
\usepackage{hyperref}

\theoremstyle{definition}
\newtheorem{problem}{Problem}

\begin{document}

\title{\vspace{-1.5cm}\textbf{Regional Quantitative Problem-Solving Circle} \\ \Large Problem Set 2: Discrete Probability \& Linearity of Expectation (Weeks 3--4)}
\author{\textbf{Lead Architect:} Mohamed Jad Srifi}
\date{\textbf{Term:} Fall 2025 -- Spring 2026}
\maketitle

\noindent\textbf{Instructions:} Complex probability and combinatorics problems often suffer from combinatorial explosion when variables are dependent. Use Indicator Random Variables and the absolute Linearity of Expectation (LoE) to bypass this dependency.

\vspace{0.5cm}

\begin{problem}[Indicator Decomposition / Circle Symmetry --- HMMT Caliber]
$N$ individuals are seated in a perfectly symmetric circle. Suddenly, each individual independently points a finger at the person immediately to their left, or the person immediately to their right, with an equal probability of $1/2$ for each direction.

Using indicator random variables, derive the expected number of individuals in the circle who are not pointed at by anyone. Explicitly state in your proof why the joint dependency of the pointing actions does not invalidate your sum.
\end{problem}

\vspace{0.5cm}

\begin{problem}[Permutation Fixed Points \& Dependent Cycle Statistics]
Let $S_n$ be the symmetric group of all permutations of the set $\{1, 2, \dots, n\}$. A permutation $\pi \in S_n$ is chosen uniformly at random.
\begin{enumerate}
    \item Let $X$ be the random variable denoting the number of fixed points in $\pi$ (i.e., the count of elements where $\pi(k) = k$). Compute $\mathbb{E}[X]$.
    \item Let $Y$ be the random variable denoting the number of transpositions (cycles of exactly length 2) in $\pi$. Compute $\mathbb{E}[Y]$.
\end{enumerate}
\end{problem}

\vspace{0.5cm}

\begin{problem}[Overcounting \& Product Space Inclusions --- AIME/PMO Caliber]
Two integers, $x$ and $y$, are drawn independently and uniformly at random from the set $\{1, 2, \dots, N\}$. 

Determine the exact probability that neither integer is the perfect square of the other (i.e., $x \neq y^2$ and $y \neq x^2$). Rigorously demonstrate your application of complementary counting and the Principle of Inclusion-Exclusion over the coordinate space.
\end{problem}

\end{document}
